\section{Introduction}

The rapid proliferation of digital technologies and online financial services has significantly increased the scale and complexity of cybercrime. In the Indian context, cybercrime incidents have grown at an alarming rate, leading to substantial economic and social impact. Recent reports indicate that cybercrime-related financial losses in India have exceeded 2.41 lakh crore INR during the period from 2020 to 2024, highlighting the urgent need for scalable and intelligent detection mechanisms.

Traditional cybercrime detection systems rely heavily on rule-based approaches or manual investigation, which are inadequate for handling the volume, diversity, and multilingual nature of modern cybercrime data. Furthermore, cybercrime incidents in India often involve multiple modalities, including textual complaints, screenshots, and fraudulent interfaces, making unimodal approaches insufficient for comprehensive analysis.

Recent advances in deep learning, particularly transformer-based architectures, have demonstrated remarkable success in natural language processing and multimodal learning. Multilingual models have enabled cross-lingual understanding, while vision-language models have improved the integration of textual and visual information. However, existing approaches are largely limited to either text-only analysis or generic multimodal tasks and do not specifically address the challenges of cybercrime detection in a multilingual and region-specific setting.

In addition, the lack of structured and annotated datasets tailored to the Indian cybercrime landscape further limits the applicability of existing models. The diversity of languages, fraud patterns, and socio-technical contexts necessitates a specialized approach that can generalize across modalities and linguistic variations while maintaining interpretability.

To address these challenges, this paper proposes a transformer-based multilingual multimodal framework for automated cybercrime detection in the Indian context. The proposed approach integrates textual and visual inputs using a cross-modal attention mechanism and leverages a curated dataset comprising 25,000 annotated samples across five Indian languages. Furthermore, a hierarchical taxonomy of 50 cybercrime categories is introduced to better capture the diversity of cybercrime incidents.

The main contributions of this work are as follows:
\begin{itemize}
    \item A novel multilingual multimodal framework for cybercrime detection using transformer-based fusion.
    \item A large-scale, annotated dataset of cybercrime incidents across multiple Indian languages.
    \item A hierarchical and clustered taxonomy comprising 50 cybercrime categories.
    \item Integration of explainable AI techniques to enhance model interpretability.
\end{itemize}